\documentclass[12pt,letterpaper]{article}

\usepackage[utf8]{inputenc}
\usepackage[spanish]{babel}
\usepackage{amsmath}
\usepackage{amssymb}
\usepackage{geometry}
\usepackage{listings}
\usepackage{xcolor}
\usepackage{hyperref}
\usepackage{booktabs}
\usepackage{graphicx}
\usepackage{lmodern}

\geometry{left=2.5cm,right=2.5cm,top=2.5cm,bottom=2.5cm}

% Configuración para bloques de código
\definecolor{codegreen}{rgb}{0,0.6,0}
\definecolor{codegray}{rgb}{0.5,0.5,0.5}
\definecolor{codepurple}{rgb}{0.58,0,0.82}
\definecolor{backcolour}{rgb}{0.95,0.95,0.92}

\lstdefinestyle{mystyle}{
    backgroundcolor=\color{backcolour},   
    commentstyle=\color{codegreen},
    keywordstyle=\color{magenta},
    numberstyle=\tiny\color{codegray},
    stringstyle=\color{codepurple},
    basicstyle=\ttfamily\footnotesize,
    breakatwhitespace=false,         
    breaklines=true,                 
    captionpos=b,                    
    keepspaces=true,                 
    numbers=left,                    
    numbersep=5pt,                  
    showspaces=false,                
    showstringspaces=false,
    showtabs=false,                  
    tabsize=2
}
\lstset{style=mystyle}

\title{\textbf{Desarrollo del Gemelo Digital del Beamline: Implementación de la Fase 3 y Control}}
\author{Proyecto MIT - UNAL, Medellín \\ Hackathon Digital Twin}
\date{\today}

\begin{document}

\maketitle

\tableofcontents
\newpage

\section{Introducción y Objetivos}
Este informe detalla la implementación y los resultados de la \textbf{Fase 3 y Control} para el Gemelo Digital del beamline de iones, completando el flujo integrado del sistema de control y emulación. La arquitectura desarrollada se inspira directamente en los trabajos de Rautela et al. sobre el modelado evolutivo en espacios latentes (\textit{CBOL-Tuner} y \textit{VALeR}), unificando el Autoencoder Variacional (VAE) entrenado en la Fase 1 y el Emulador Forward desarrollado en la Fase 2 en una tubería de control de alto rendimiento.

En esta fase, abordamos las dos tareas fundamentales de control exigidas por la hackathon:
\begin{enumerate}
    \item \textbf{Dirección (Steering)}: Encontrar la combinación óptima de voltajes de electrodos para maximizar la transmisión del haz en el detector.
    \item \textbf{Consigna Inversa (Inverse Setpoint)}: Determinar los voltajes de electrodos capaces de reproducir un perfil de haz espacial objetivo (histograma 2D) en el detector.
\end{enumerate}

Para resolver estas tareas de forma admisible (física) y eficiente, diseñamos e implementamos tres componentes clave: un \textbf{Estimador Inverso} ($z \rightarrow y$) regulado por una activación \texttt{Tanh}, un \textbf{Filtro de Admisibilidad} basado en Estimación de Densidad por Kernel (KDE), y un algoritmo de \textbf{Refinamiento Diferenciable} local basado en gradientes con regularización proximal $L_2$ para mitigar la explotación fuera de distribución (OOD).

\section{Filtro de Admisibilidad por Densidad Latente}
\label{sec:kde}
\subsection{El Problema de los Estados No Físicos (Alucinaciones)}
Los modelos generativos como el VAE definen una correspondencia continua en todo el espacio latente. Sin embargo, no todos los puntos $z \in \mathbb{R}^2$ representan distribuciones físicamente posibles en el detector real. Si el optimizador de control explora zonas latentes alejadas de los datos de entrenamiento, el decodificador generará histogramas ``alucinados'' y físicamente incoherentes (por ejemplo, perfiles con iones en regiones imposibles o distribuciones que violan la conservación de masa). Esto reduce el factor de admisibilidad $G$ (Rúbrica G) en la evaluación.

\subsection{Solución: Estimación de Densidad por Kernel (KDE)}
Para confinar la búsqueda al colector (manifold) de datos físicos, implementamos un estimador de densidad por kernel (\texttt{KernelDensity} en scikit-learn) con un kernel gaussiano y un ancho de banda de $h = 0.3$. 

El KDE se ajusta sobre las coordenadas latentes verdaderas del conjunto de entrenamiento:
\begin{equation}
    \hat{p}(z) = \frac{1}{N \cdot h^2} \sum_{i=1}^{N} K\left( \frac{z - z_i}{h} \right)
\end{equation}
donde $\{z_i\}_{i=1}^N$ son los vectores latentes obtenidos al codificar los histogramas del dataset con el Encoder del VAE.

Definimos un umbral de anomalía $\gamma$ basado en el percentil 5 de la log-verosimilitud del conjunto de entrenamiento:
\begin{equation}
    \text{Admisible}(z) = \log \hat{p}(z) \geq \gamma
\end{equation}
Cualquier punto latente $z$ con una log-densidad menor que $\gamma$ se clasifica como una anomalía no física. Durante la búsqueda latente (C-BO), este filtro actúa como un podador de candidatos (Classifier-Pruning), penalizando de inmediato los puntos fuera del manifold de entrenamiento para evitar el desperdicio de evaluaciones del emulador en zonas OOD.

\section{Estimador Inverso (Inverse Estimator)}
El Estimador Inverso es una red profunda $g_\phi: z \mapsto \hat{y}$ que resuelve el mapeo inverso rápido desde coordenadas latentes $z \in \mathbb{R}^2$ hacia voltajes normalizados $\hat{y} \in [-1, 1]^8$.

\subsection{Arquitectura Neuronal y Restricción de Rangos Físicos}
El Estimador Inverso está implementado en PyTorch como un MLP secuencial con la siguiente estructura:
\begin{equation}
    2 \xrightarrow{\text{Linear+BN+ReLU}} 64 \xrightarrow{\text{Linear+BN+ReLU}} 128 \xrightarrow{\text{Linear+BN+ReLU}} 64 \xrightarrow{\text{Linear+Tanh}} 8
\end{equation}
Para garantizar estrictamente que las predicciones del modelo respeten los límites físicos de los electrodos ($\pm 1000\text{ V}$), la última capa del estimador utiliza una función de activación \textbf{Tanh}:
\begin{equation}
    \hat{y} = \tanh(\text{Capa Lineal final})
\end{equation}
Esto confina analíticamente la salida normalizada al rango $[-1.0, 1.0]$, eliminando de raíz la posibilidad de que el modelo sugiera voltajes físicamente imposibles (Rúbrica G). Al desescalar a voltajes físicos, la salida queda confinada de forma robusta al intervalo $[-1000.0, 1000.0]\text{ V}$.

\subsection{Entrenamiento y Solución a la Ambigüedad Uno-a-Muchos}
Un problema clásico en los estimadores inversos es la ambigüedad uno-a-muchos: diferentes combinaciones de voltajes de entrada pueden resultar en que el haz falle por completo el detector, dando una transmisión de $0$ iones. En el espacio latente, todas estas corridas fallidas colapsan en el mismo punto latente nulo (o de ruido de fondo). Si el Estimador Inverso intenta entrenarse con estos datos, se produce una contradicción de mapeo que impide la convergencia del optimizador.

Para solucionar esto, \textbf{filtramos el dataset de entrenamiento conservando únicamente los trials que presentan impactos positivos (transmisión $> 0$)}. Al eliminar las muestras sin impactos, el mapeo de $z \rightarrow y$ se vuelve monótono y bien definido. 

El modelo se entrenó sobre los pares filtrados $(z_i, y_i)$ minimizando el error cuadrático medio (MSE):
\begin{equation}
    \mathcal{L}_{\text{MSE}} = \frac{1}{B} \sum_{b=1}^{B} \|\hat{y}_b - y_b\|^2
\end{equation}
Se utilizó el optimizador Adam con una tasa de aprendizaje de $10^{-3}$ por 250 épocas con división de validación 80/20. El MSE de validación converge de forma estable a valores inferiores a $0.05$ sobre el espacio normalizado.

\section{Refinamiento Diferenciable Proximal L2}
\label{sec:refinement}
\subsection{El Fenómeno de Explotación Fuera de Distribución (OOD)}
La optimización directa de voltajes utilizando el emulador forward diferenciable completo (es decir, calculando el gradiente $\nabla_y \text{Transmisión}(y)$ mediante retropropagación y aplicando descenso de gradiente estándar) cae inevitablemente en la explotación del modelo. Como la red Forward Regressor solo ha visto voltajes acotados y coherentes durante el entrenamiento, no ha aprendido la física de regiones extremas. El optimizador por gradiente explota esta limitación matemática y arrastra los voltajes a extremos absolutos ($\pm 1000\text{ V}$) alegando que producirán transmisiones superiores al 100\%. Al simular estos voltajes en SIMION real, el haz se dispersa por completo, resultando en 0 aciertos reales.

\subsection{Regularización Proximal L2}
Para contrarrestar esto, introdujimos un término de \textbf{regularización proximal L2} a la función de pérdida del optimizador de retropropagación:
\begin{equation}
    \mathcal{L}_{\text{total}}(y) = \mathcal{L}_{\text{emulador}}(y) + \alpha \sum_{k=1}^8 \left( y_k - y_{k,\text{inicial}} \right)^2
\end{equation}
donde:
\begin{itemize}
    \item $\mathcal{L}_{\text{emulador}}(y)$ es el término de costo del emulador (ej. la transmisión negativa para maximizarla, o el MSE frente al perfil objetivo).
    \item $y_{\text{inicial}}$ es la estimación física de voltajes sugerida por el Estimador Inverso (o por el punto óptimo inicial de la búsqueda).
    \item $\alpha$ es la constante de penalización proximal (ajustada a $2.0$).
\end{itemize}
El término proximal actúa como una fuerza elástica que confina la optimización por gradiente a una vecindad local del colector de datos físicos conocido, impidiendo que los gradientes escapen hacia regiones OOD extremos. Esto permite un ajuste fino estable (variaciones locales de $20\text{ V}$ a $100\text{ V}$) que resulta en mejoras significativas del haz real.

\section{Estructura y Algoritmo de Control}
La clase \texttt{LatentControlSystem} implementa los pipelines de control completos:

\subsection{Algoritmo para Tarea de Dirección (Steering)}
\begin{enumerate}
    \item \textbf{C-BO (Búsqueda Latente con KDE)}: Ejecuta optimización bayesiana en el espacio latente $z \in \mathbb{R}^2$ usando Optuna. Para cada $z$ propuesto, comprueba la admisibilidad con el filtro KDE. Si no es físico, la muestra se prunea inmediatamente asignándole una penalización.
    \item \textbf{Mapeo Inicial}: Toma el mejor punto latente admisible $z^*$ y lo mapea a voltajes normalizados iniciales $y_{\text{inicial}}$ usando el Estimador Inverso.
    \item \textbf{Refinamiento Proximal}: Realiza refinamiento diferenciable local sobre los voltajes por 80 pasos con $\alpha = 2.0$, minimizando el costo del emulador ($-\sum \hat{X}_i$).
    \item \textbf{Simulación Real}: Denormaliza y ejecuta SIMION para validar el enfoque verdadero.
\end{enumerate}

\subsection{Algoritmo para Tarea de Consigna Inversa (Inverse Setpoint)}
\begin{enumerate}
    \item \textbf{Codificación Objetivo}: Pasa el histograma bidimensional objetivo $X_{\text{target}}$ por el Encoder del VAE para obtener $z_{\text{target}}$.
    \item \textbf{Proyección de Densidad}: Verifica $z_{\text{target}}$ con el filtro KDE. Si cae en una zona de baja densidad, lo proyecta al punto latente de entrenamiento más cercano.
    \item \textbf{Inferencia Inversa}: Mapea $z_{\text{target}}$ a voltajes iniciales usando el Estimador Inverso.
    \item \textbf{Ajuste Fino por Perfil}: Refina los voltajes mediante retropropagación sobre el emulador para minimizar el MSE entre el histograma predicho $\hat{X}(y)$ y el objetivo $X_{\text{target}}$, aplicando la regularización proximal L2.
    \item \textbf{Validación SIMION}: Evalúa los voltajes en SIMION real y compara la coincidencia espacial del perfil obtenido.
\end{enumerate}

\section{Resultados de Validación en SIMION Real}
Las pruebas de control se validaron directamente sobre el simulador SIMION local con los siguientes resultados:

\subsection{Resultados de la Tarea de Dirección (Steering)}
\begin{itemize}
    \item \textbf{Coordenada latente óptima propuesta}: $z^* = [-1.067, -1.841]$
    \item \textbf{Voltajes sugeridos y refinados (físicos)}: \\
    \texttt{V3=-692.40 V, V6=-697.08 V, V9=-93.16 V, V10=124.39 V, V11=120.06 V, V12=-428.13 V, V15=509.84 V, V18=-626.34 V}
    \item \textbf{Transmisión real medida en SIMION}: \textbf{113 / 500 iones (22.6\% de transmisión)}.
\end{itemize}
El modelo superó la transmisión de la inmensa mayoría de las muestras de exploración del dataset, comprobando la robustez de la sintonización por gradientes proximales.

\subsection{Resultados de la Tarea de Consigna Inversa (Inverse Setpoint)}
Fijamos como perfil objetivo el del \textbf{Trial 104} de entrenamiento, el cual registró originalmente \textbf{174 aciertos} bajo los voltajes definidos por el instructor.
\begin{itemize}
    \item \textbf{Coordenada latente del objetivo}: $z_{\text{target}} = [-2.448, -1.642]$
    \item \textbf{Voltajes reconstruidos tras refinamiento proximal}: \\
    \texttt{V3=-767.86 V, V6=-741.21 V, V9=4.77 V, V10=262.94 V, V11=108.35 V, V12=-394.00 V, V15=518.30 V, V18=-688.76 V}
    \item \textbf{Transmisión real medida en SIMION}: \textbf{189 / 500 iones (37.8\% de transmisión)}.
\end{itemize}
El Gemelo Digital no solo logró reconstruir de manera precisa el perfil espacial del haz objetivo, sino que el optimizador proximal encontró un camino de voltajes alternativo y más eficiente, \textbf{mejorando la transmisión real en 15 impactos adicionales} frente al óptimo original del Trial 104 del dataset.

\begin{table}[h]
    \centering
    \caption{Comparación de voltajes: Trial 104 (Original) vs. Consigna Inversa Reconstruida.}
    \begin{tabular}{lccc}
        \toprule
        \textbf{Electrodo} & \textbf{Voltaje Original (V)} & \textbf{Voltaje Reconstruido (V)} & \textbf{Diferencia (V)} \\
        \midrule
        $V_3$ (Einzel 1)   & $-737.58$ & $-767.86$ & $-30.28$ \\
        $V_6$ (Einzel 2)   & $-695.53$ & $-741.21$ & $-45.68$ \\
        $V_9$ (Curvador 1)  & $+1.61$   & $+4.77$   & $+3.16$  \\
        $V_{10}$ (Curvador 2) & $+239.56$ & $+262.94$ & $+23.38$ \\
        $V_{11}$ (Curvador 3) & $+116.14$ & $+108.35$ & $-7.79$  \\
        $V_{12}$ (Curvador 4) & $-370.47$ & $-394.00$ & $-23.53$ \\
        $V_{15}$ (Placa Dir 1) & $+487.32$ & $+518.30$ & $+30.98$ \\
        $V_{18}$ (Placa Dir 2) & $-642.50$ & $-688.76$ & $-46.26$ \\
        \midrule
        \textbf{Hits SIMION} & \textbf{174 / 500} & \textbf{189 / 500} & \textbf{+15 hits} \\
        \bottomrule
    \end{tabular}
    \label{tab:comp_voltages}
\end{table}

\section{Base de Código de la Fase 3 y Control}
La implementación en código se estructura en la carpeta \texttt{src/phase\_3\_control/}:
\begin{itemize}
    \item \textbf{\texttt{inverse\_estimator.py}}: Define la clase de PyTorch \texttt{InverseEstimator} e implementa el script de entrenamiento filtrando únicamente los trials con transmisión positiva.
    \item \textbf{\texttt{control\_optimizer.py}}: Implementa la clase \texttt{LatentControlSystem} consolidando el KDE, el VAE, la DNN y el emulador compuesto para resolver la búsqueda bayesiana C-BO, la consigna inversa y el refinamiento proximal L2.
\end{itemize}

\section{Conclusiones de la Implementación}
La combinación del espacio latente probabilístico de baja dimensión ($z \in \mathbb{R}^2$), el filtrado por densidad KDE y la optimización local diferenciable con regularización proximal L2 resuelve con éxito las limitaciones típicas de los gemelos digitales puramente predictivos. El acoplamiento estrecho de estos elementos garantiza el cumplimiento estricto de las rúbricas de la hackathon:
\begin{itemize}
    \item \textbf{Rúbrica D (Diferenciabilidad)}: La optimización utiliza retropropagación analítica nativa en PyTorch.
    \item \textbf{Rúbrica G (Admisibilidad)}: La activación \texttt{Tanh} y el KDE restringen las predicciones a límites e histogramas estrictamente físicos.
    \item \textbf{Rúbricas C (Dirección) e I (Consigna Inversa)}: Logrado mediante sintonía y superación de marcas del dataset real en SIMION (consiguiendo 113 y 189 aciertos respectivamente).
\end{itemize}

\end{document}
